\section{Software Product Definition}
\begin{longtable}[ht]{|L{3.5cm}|L{5.5cm}|L{5.5cm}|}\hline%
  \rowcolor{fkblue}%
  \fkHdrRow{Item} & \fkHdrRow{Number} & \fkHdrRow{Title/Name}\ER%
  \endhead%
  Software Architecture&%
  \swarcnum&%
  \swarctitle\ER%
  %
  Software Failure Modes and Effects Analysis&%
  \swfmeanum&%
  \swfmeatitle\ER%
  %
  Software of Unknown Provenance Risk Analysis&%
  \soupnum&%
  \souptitle\ER%
  %
  Software Requirements Specifications&%
  \swreqmntnums&%
  \swreqmnttitles\ER%
  %
  Systems Requirements Specifications&%
  \sysreqmntnums&%
  \sysreqmnttitles\ER%
  %
  \caption{Software Product Definition References}%
  \label{table:2}%
\end{longtable}%

\tlcVspace

\textbf{NOTE:} As the software product is software only, its SSR is written with
sufficient detail to represent the Software Requirement Specification for the
purposes of Software Product Definition. The \Gls{soup} Requirements
Specification covers the gap between the SSR and the standard Software
Requirements Specification.

\tlcVspace

The deliverables listed in Table 2 identify significant design inputs that
define the software product. As necessary, these deliverables will be updated as
development of the software product proceeds. For a comprehensive enumeration of
all design inputs for the software product, see the \gls{company} design history
file.

\tlcVspace

For Software Safety classification / Level of Concern, see the \Gls{rmp} in
\gls{company} \Gls{dhf}.

\tlcVspace
Software will follow standards from the regulatory assessment in \gls{company}
\Gls{dhf}.
